\section{\label{sec:hyperchords} Plane-dihedral construction in three dimensions}

We now promote the planar half-angle construction of Section~\ref{sec:planar} to three dimensions,
obtaining the central geometric objects of this paper: two offset circle sections on a sphere,
which we call the \emph{plane-dihedral faces}.

\subsection{Setup: rotated equatorial plane and the highest equator point}
\label{sec:hyperchords_setup}
\paragraph{Geometric scaling convention.}
A two-component $SU(2)$ spinor $\ssp\in\mathbb{C}^2$ carries a natural quadratic invariant
$u:=\ssp^\dagger\ssp$.
In this work we use $u$ as the \emph{geometric scale parameter}.
Concretely, rather than embedding $\ssp$ directly, we embed the rescaled spinor
\begin{equation}
\label{eq:scaled_spinor}
\widetilde{\ssp} := \|\ssp\|\,\ssp = \sqrt{u}\,\ssp,
\end{equation}
whose Euclidean norm is
\begin{equation}
\label{eq:scaled_norm}
\|\widetilde{\ssp}\| = \sqrt{\widetilde{\ssp}^\dagger\widetilde{\ssp}}
= \sqrt{(\sqrt{u}\,\ssp)^\dagger(\sqrt{u}\,\ssp)} = \sqrt{u(\ssp^\dagger\ssp)} = u.
\end{equation}
This is the reason the sphere used in the visualization has radius $u$.
In dimensional terms, if one assigns spinors the conventional length-dimension $[\ssp]=L^{1/2}$ (so that quadratic spinor bilinears carry dimension $L$), then $u=\ssp^\dagger\ssp$ has units of length and $S_u$ is a genuine length-scale object.
For normalized quantum states one has $u=1$, and the construction reduces to the unit-radius case.
For notational simplicity, we continue to write $\ssp$ for the input spinor and use $u:=\ssp^\dagger\ssp$ to set the scene scale.

Fix the sphere of radius $u$,
\[
S_u := \{\,\mathbf{x}\in\Rt : \norm{\mathbf{x}}=u\,\},
\]
with $u:=\ssp^\dagger\ssp$ as in \eqref{eq:scaled_spinor}--\eqref{eq:scaled_norm}. Note that the geometry of the two faces depends only on $(\theta,\phi)$, while $u$ acts as an overall similarity scale and $\gamma$ acts as a common in-face marker rotation.


and the reference equatorial plane $E:=\xyplane$ with equator $e:=E\cap S_u$.
The primary datum is a rotated equatorial plane $E'=\RotMat E$, for some $\RotMat\in\sot$.
Unprimed symbols refer to the reference configuration; primed symbols to the rotated one.

Let $\vv$ be a chosen unit normal to $E'$ (so $E'$ determines $\vv$ up to sign), and define
\[
\theta\in[0,\pi]
\qquad\text{by}\qquad
\vv\cdot \eez = \cos\theta .
\]
Equivalently, $\theta$ is the dihedral angle between $E'$ and $E$ (since $\eez$ is the unit normal to $E$).

Let
\[
e' := E'\cap S_u
\]
denote the rotated equator. We single out a canonical point on $e'$ by taking the maximizer of the
$z$-coordinate:
\[
Q \;:=\; \arg\max_{\mathbf{x}\in e'} \;(\mathbf{x}\cdot\eez).
\]
For $\theta\in(0,\pi)$ this point is unique and can be written explicitly.
Indeed, the unit vector in $E'$ that maximizes $\mathbf{x}\cdot\eez$ is the normalized projection of $\eez$
onto $E'$, i.e.
\[
\widehat{\mathbf{h}}
:=
\frac{\eez-(\eez\cdot\vv)\,\vv}{\norm{\eez-(\eez\cdot\vv)\,\vv}}
=
\frac{\eez-\cos\theta\,\vv}{\sin\theta}
\in E',
\qquad
Q := u\,\widehat{\mathbf{h}}.
\]
In particular,
\[
Q\cdot\eez = u\,\sin\theta,
\qquad
\text{and}\qquad
\vv\cdot Q = 0.
\]
(For the degenerate cases $\theta=0$ or $\theta=\pi$, the entire equator $e'$ lies in the $xy$-plane and
the maximizer is not unique; in that case we fix $Q$ by choosing an azimuth $\phi$ as discussed in
Section~\ref{sec:hyperchords_equator}.)

\subsection{Definition of the two faces via dihedral bisectors}
\label{sec:hyperchords_faces}

The two face planes $U'$ and $L'$ are defined as follows.

First form the two unit \emph{bisector normals}
\[
\nn_\uparrow := \frac{\eez+\vv}{\norm{\eez+\vv}}
= \frac{\eez+\vv}{2\cos(\theta/2)},
\qquad
\nn_\downarrow := \frac{\eez-\vv}{\norm{\eez-\vv}}
= \frac{\eez-\vv}{2\sin(\theta/2)},
\]
where the half-angle normalizations use $\norm{\eez+\vv}^2 = 2+2\cos\theta = 4\cos^2(\theta/2)$ and
$\norm{\eez-\vv}^2 = 2-2\cos\theta = 4\sin^2(\theta/2)$.
A direct computation gives
\[
\nn_\uparrow\cdot\nn_\downarrow
\;\propto\;
(\eez+\vv)\cdot(\eez-\vv)=\eez\cdot\eez-\vv\cdot\vv=0,
\]
so the bisector normals are orthogonal.

Now define the \emph{offset} planes as the unique translates of the corresponding bisector planes
through the origin that pass through $Q$:
\[
U' := \{\,\mathbf{x}\in\Rt : \nn_\uparrow\cdot \mathbf{x} = \nn_\uparrow\cdot Q\,\},
\qquad
L' := \{\,\mathbf{x}\in\Rt : \nn_\downarrow\cdot \mathbf{x} = \nn_\downarrow\cdot Q\,\}.
\]
By construction $Q\in U'\cap L'$ and $U'\perp L'$ (since their normals are orthogonal).

The two plane-dihedral circles are then defined as the corresponding sphere sections
\[
c_\uparrow := U'\cap S_u,
\qquad
c_\downarrow := L'\cap S_u.
\]

\paragraph{Offsets and radii (half-angle law).}
Because $Q\in E'$ we have $\vv\cdot Q=0$, hence
\[
\nn_\uparrow\cdot Q
=
\frac{(\eez+\vv)\cdot Q}{2\cos(\theta/2)}
=
\frac{\eez\cdot Q}{2\cos(\theta/2)}
=
\frac{u\sin\theta}{2\cos(\theta/2)}
=
u\sin(\theta/2),
\]
and similarly
\[
\nn_\downarrow\cdot Q
=
\frac{(\eez-\vv)\cdot Q}{2\sin(\theta/2)}
=
\frac{\eez\cdot Q}{2\sin(\theta/2)}
=
\frac{u\sin\theta}{2\sin(\theta/2)}
=
u\cos(\theta/2).
\]
Therefore the signed distances of the planes from the origin are
\[
\mathrm{dist}(U',0)=u\sin(\theta/2),
\qquad
\mathrm{dist}(L',0)=u\cos(\theta/2).
\]
A plane at distance $d$ from the origin cuts $S_u$ in a circle of radius $\sqrt{u^2-d^2}$, hence
\[
r_\uparrow = \sqrt{u^2-u^2\sin^2(\theta/2)} = u\cos(\theta/2),
\qquad
r_\downarrow = \sqrt{u^2-u^2\cos^2(\theta/2)} = u\sin(\theta/2),
\]
which is precisely the desired half-angle scaling.

\subsection{Connection to the equator and definition of the azimuth \texorpdfstring{$\phi$}{phi}}
\label{sec:hyperchords_equator}

Each plane-dihedral circle is tangent to the reference equator $e$ at a single point, and these two tangent
points are antipodal.

To see this, intersect $U'$ with the $xy$-plane $E$ (i.e. set $z=0$). Using
$\nn_\uparrow=(\eez+\vv)/(2\cos(\theta/2))$ and $\nn_\uparrow\cdot Q=u\sin(\theta/2)$, the plane equation
$\nn_\uparrow\cdot\mathbf{x}=\nn_\uparrow\cdot Q$ reduces on $E$ to
\[
\vv\cdot\mathbf{x} = u\sin\theta
\qquad (z=0).
\]
Writing $\vv=(v_x,v_y,v_z)$ with $v_z=\cos\theta$ and $\sqrt{v_x^2+v_y^2}=\sin\theta$, this is a line in $E$
at distance
\[
\frac{u\sin\theta}{\sqrt{v_x^2+v_y^2}} = \frac{u\sin\theta}{\sin\theta}=u
\]
from the origin. Hence the line $\ell_\uparrow := U'\cap E$ is tangent to the equator circle
$e=\{(x,y,0):x^2+y^2=u^2\}$ at a \emph{unique} point, which we denote by $P_\uparrow$.
Explicitly, let $\widehat{\vv}_{xy}$ be the unit projection of $\vv$ onto $E$,
\[
\widehat{\vv}_{xy} := \frac{1}{\sin\theta}\,(v_x,v_y,0),
\]
then the tangency point is
\[
P_\uparrow := u\,\widehat{\vv}_{xy}\in e.
\]
The same calculation for $L'$ gives a tangent line $\ell_\downarrow := L'\cap E$ whose unique tangency point
is the antipode
\[
P_\downarrow := -P_\uparrow.
\]

We define the azimuthal angle $\phi$ by writing
\[
P_\uparrow = u(\cos\phi,\sin\phi,0),
\qquad
P_\downarrow = u(\cos(\phi+\pi),\sin(\phi+\pi),0).
\]
In the visualization developed later, this $\phi$ is the geometric carrier of the \emph{relative phase}:
rotating the entire configuration about the $z$-axis rigidly shifts $\phi$ by the same amount, while the
half-angle structure of the faces will induce the characteristic spinorial response.

\begin{figure}[ht!]
    \centering
    \includegraphics[width=0.7\columnwidth]{viz_figures/hyperchords_3d_py.png}
    \caption{Plane-dihedral construction. The rotated equatorial plane $E'$ (primary datum) fixes the unit normal
    $\vv$ and the dihedral angle $\theta$ with the reference plane $E=\xyplane$. The rotated equator
    $e'=E'\cap S_u$ has a distinguished highest point $Q$. The two face planes $U'$ and $L'$ are the
    unique translates (through $Q$) of the two dihedral-bisector planes between $E$ and $E'$, with normals
    proportional to $\eez\pm\vv$. Their sphere sections are the circles $c_\uparrow$ (teal/bluish-green, thicker line, equator contact marked with $\blacktriangle$) and $c_\downarrow$ (orange, thinner line, equator contact marked with $\blacktriangledown$), with
    radii $r_\uparrow=u\cos(\theta/2)$ and $r_\downarrow=u\sin(\theta/2)$. Each circle is tangent to the
    reference equator $e$ at a unique point: $P_\uparrow$ for $c_\uparrow$ and $P_\downarrow=-P_\uparrow$
    for $c_\downarrow$. The azimuth of $P_\uparrow$ is $\phi$.}
    \label{fig:hyperchords_schematic}
\end{figure}

\subsection{Geometric meaning of the phases}
\label{sec:hyperchords_phase}

Once the spinor is written in the symmetric half-angle phase split used throughout this work,
the two faces carry distinguished in-face diameter directions whose \emph{relative} rotation encodes
the relative phase $\phi$, while their \emph{common} rotation encodes the fiber (global) phase $\gamma$.
(We make this explicit after introducing the $SU(2)$ action on spinors in Section~\ref{sec:su2_viz}, using the symmetric phase split
$\psi_\uparrow=\gamma-\phi/2$, $\psi_\downarrow=\gamma+\phi/2$.)